% --- IMPORTS ---
\documentclass[leqno]{article}
\usepackage{graphicx}
\usepackage{dsfont}
\usepackage{mathtools}
\usepackage{amssymb}
\usepackage{amsmath}
\usepackage{amsthm}
\usepackage{float}
\usepackage{ragged2e}
\usepackage{tensor}
\usepackage{fancyhdr}
\usepackage{empheq}
\usepackage[most]{tcolorbox}
\usepackage{enumitem}
\usepackage{dirtytalk}
\usepackage{marginnote}
\usepackage{microtype}
\usepackage[
  a4paper,
  left=0.8in,
  right=1in,
  top=1in,
  bottom=0.5in,
  headheight=14pt,
  headsep=0.4in
]{geometry}
\setlength{\parindent}{0cm}
% --- TikZ Packages ---
\usepackage{pgfplots}
\pgfplotsset{compat=1.18}
\usepgfplotslibrary{fillbetween, polar}
\usetikzlibrary{calc, positioning}
\usetikzlibrary{angles, quotes}
\usetikzlibrary{arrows.meta}
\usetikzlibrary{decorations.pathreplacing, decorations.markings}
\usetikzlibrary{patterns}
\usetikzlibrary{intersections}
% --- URL Packages ---
\usepackage[colorlinks=true,linkcolor=red,citecolor=magenta,urlcolor=black]{hyperref}%
\urlstyle{same}
% --- END OF IMPORTS ---

% --- CUSTOM COMMANDS ---
\RenewDocumentCommand{\marks}{o m}{%
  \IfValueTF{#1}%
    {\marginnote{\raggedright\textbf{[#2]}}[#1]}%
    {\marginnote{\raggedright\textbf{[#2]}}}%
}
\newcommand{\vspacerule}[2]{%
  \par\vspace*{#1}%
  \noindent\hrulefill\par
  \vspace*{#2}%
}
\definecolor{scarlet}{RGB}{205, 40, 75}
\newcommand{\questionitem}{%
  \item\phantomsection
  \addcontentsline{toc}{section}{Question \number\value{enumi}}%
}
\renewcommand{\contentsname}{Questions}
% --- END OF CUSTOM COMMANDS ---

% --- HEADER CONFIG ---
\pagestyle{fancy}
\fancyhead{}
\fancyfoot{}
\renewcommand{\headrulewidth}{0.9pt}
\lhead{\textit{Solutions} - Integration with Inverse Trig}
\chead{}
\rhead{\href{https://danieljsmith.org}{danieljsmith.org}}
% --- END OF HEADER CONFIG ---

\begin{document}
\vspace*{20pt}
\tableofcontents
\newpage
\begin{enumerate}[
  label=\textbf{\arabic*.},
  leftmargin=*,
  topsep=0pt,
  partopsep=0pt,
  parsep=0pt
]

% ---- Question 1 ----
\questionitem
% [Source: _QBT___Solns__Integration_with_Inverse_Trig, Q1]

Find the exact value of the improper integral
\[
\int_0^4 \frac{1}{\sqrt{x(4-x)}}\,\mathrm{d}x
\]
\marks[-30pt]{5}

\vspace*{10pt}

\begin{tcolorbox}[colback=white, colframe=scarlet, title={\textbf{Solution}}, breakable, grow to left by=6mm, grow to right by=10mm]
Let
\[
I=\int_0^4 \frac{1}{\sqrt{x(4-x)}}\,\mathrm{d}x
\]

This is an improper integral because the denominator is \(0\) at \(x=0\) and \(x=4\). So we write
\[
I=\lim_{\substack{a\to 0^+\\ b\to 4^-}}\int_a^b \frac{1}{\sqrt{x(4-x)}}\,\mathrm{d}x
\]

Now simplify the expression inside the square root:
\[
x(4-x)=4x-x^2=4-(x^2-4x+4)=4-(x-2)^2
\]

So
\[
I=\lim_{\substack{a\to 0^+\\ b\to 4^-}}\int_a^b \frac{1}{\sqrt{4-(x-2)^2}}\,\mathrm{d}x
\]

Let
\[
u=x-2
\]
Then
\[
\mathrm{d}u=\mathrm{d}x
\]
and the limits become
\[
x=a \Rightarrow u=a-2,\qquad x=b \Rightarrow u=b-2
\]

Hence
\[
I=\lim_{\substack{a\to 0^+\\ b\to 4^-}}\int_{a-2}^{b-2} \frac{1}{\sqrt{4-u^2}}\,\mathrm{d}u
\]

Using the standard result
\[
\int \frac{1}{\sqrt{a^2-u^2}}\,\mathrm{d}u=\arcsin \!\left(\frac{u}{a}\right)+c
\]
with \(a=2\), we get
\[
I=\lim_{\substack{a\to 0^+\\ b\to 4^-}}\left[\arcsin \!\left(\frac{u}{2}\right)\right]_{a-2}^{b-2}
\]

So
\[
I=\lim_{\substack{a\to 0^+\\ b\to 4^-}}\left(\arcsin \!\left(\frac{b-2}{2}\right)-\arcsin \!\left(\frac{a-2}{2}\right)\right)
\]

Now
\[
\frac{b-2}{2}\to 1 \quad \text{as } b\to 4^-
\]
and
\[
\frac{a-2}{2}\to -1 \quad \text{as } a\to 0^+
\]

Therefore
\[
I=\arcsin (1)-\arcsin (-1)=\frac{\pi}{2}-\left(-\frac{\pi}{2}\right)=\pi
\]

The exact value of the integral is \(\pi\)
\end{tcolorbox}


\newpage

% ---- Question 2 ----
\questionitem
% [Source: _QBT___Solns__Integration_with_Inverse_Trig, Q2]

Show that
\[
\int_{0}^{3} \frac{1}{x^2 - 2x + 5}\,\mathrm{d}x = \frac{1}{2}\arctan (3)
\]
\marks[-20pt]{5}

\vspace*{10pt}

\begin{tcolorbox}[colback=white, colframe=scarlet, title={\textbf{Solution}}, breakable, grow to left by=6mm, grow to right by=10mm]
Let
\[
I=\int_{0}^{3}\frac{1}{x^2-2x+5}\,\mathrm{d}x
\]

First complete the square in the denominator:
\[
x^2-2x+5=(x-1)^2+4=(x-1)^2+2^2
\]

So
\[
I=\int_{0}^{3}\frac{1}{(x-1)^2+2^2}\,\mathrm{d}x
\]

Now use the standard result from the formula booklet:
\[
\int \frac{1}{u^2+a^2}\,\mathrm{d}u=\frac{1}{a}\arctan \!\left(\frac{u}{a}\right)+c
\]

Here, take \(u=x-1\), so \(\mathrm{d}u=\mathrm{d}x\), and \(a=2\). Hence
\[
I=\frac{1}{2}\left[\arctan \!\left(\frac{x-1}{2}\right)\right]_{0}^{3}
\]

Substituting the limits gives
\begin{align*}
I&=\frac{1}{2}\left(\arctan \!\left(\frac{3-1}{2}\right)-\arctan \!\left(\frac{0-1}{2}\right)\right)\\
&=\frac{1}{2}\left(\arctan (1)-\arctan \!\left(-\frac12\right)\right)\\
&=\frac{1}{2}\left(\frac{\pi}{4}+\arctan \!\left(\frac12\right)\right)
\end{align*}

To simplify this, let
\[
\theta=\frac{\pi}{4}+\arctan \!\left(\frac12\right)
\]

Then
\begin{align*}
\tan\theta
&=\frac{\tan\left(\frac{\pi}{4}\right)+\tan\left(\arctan \!\left(\frac12\right)\right)}
{1-\tan\left(\frac{\pi}{4}\right)\tan\left(\arctan \!\left(\frac12\right)\right)}\\
&=\frac{1+\frac12}{1-1\cdot\frac12}\\
&=\frac{\frac32}{\frac12}\\
&=3
\end{align*}

Also, both \(\frac{\pi}{4}\) and \(\arctan \!\left(\frac12\right)\) are between \(0\) and \(\frac{\pi}{2}\), so \(\theta\) is in the first quadrant. Therefore
\[
\theta=\arctan (3)
\]

Hence
\[
I=\frac{1}{2}\arctan (3)
\]

So
\[
\int_{0}^{3}\frac{1}{x^2-2x+5}\,\mathrm{d}x=\frac{1}{2}\arctan (3)
\]

as required.
\end{tcolorbox}


\newpage


% ---- Question 3 ----
\questionitem
% [Source: _QBT___Solns__Integration_with_Inverse_Trig, Q3]

\begin{enumerate}[label=(\roman*)]
  \item Prove, from definitions involving exponentials, that
  \[
  1-\tanh^2u=\operatorname{sech}^2u
  \]
  \marks[-20pt]{3}
  \item Prove that, for $|y|<1$,
  \[
  \operatorname{artanh}y=\frac{1}{2}\ln\left(\frac{1+y}{1-y}\right)
  \]
  \marks[-20pt]{4}
  \item Use the substitution $x=4\tanh u$ to show that, for $|x|<4$,
  \[
  \int \frac{1}{16-x^2}\,\mathrm{d}x=\frac{1}{4}\operatorname{artanh}\frac{x}{4}+c
  \]
  Hence deduce that
  \[
  \int \frac{1}{16-x^2}\,\mathrm{d}x=\frac{1}{8}\ln\left(\frac{4+x}{4-x}\right)+c
  \]
  where $c$ is an arbitrary constant. \marks{6} \\
  \item By first expressing $-t^2+6t+7$ in completed square form, show that
  \[
  \int_{1}^{5} \frac{1}{-t^2+6t+7}\,\mathrm{d}t=\frac{1}{4}\ln 3
  \]
  \marks[-20pt]{5}
\end{enumerate}

\vspace*{10pt}

\begin{tcolorbox}[colback=white, colframe=scarlet, title={\textbf{Solution}}, breakable, grow to left by=6mm, grow to right by=10mm]
\begin{enumerate}[label=\textbf{(\roman*)}]
\item Using the exponential definitions,
\[
\tanh u=\frac{e^u-e^{-u}}{e^u+e^{-u}}
\qquad\text{and}\qquad
\operatorname{sech}u=\frac{2}{e^u+e^{-u}}
\]

Then
\[
\begin{aligned}
1-\tanh^2u
&=1-\left(\frac{e^u-e^{-u}}{e^u+e^{-u}}\right)^2 \\
&=\frac{(e^u+e^{-u})^2-(e^u-e^{-u})^2}{(e^u+e^{-u})^2} \\
&=\frac{(e^{2u}+2+e^{-2u})-(e^{2u}-2+e^{-2u})}{(e^u+e^{-u})^2} \\
&=\frac{4}{(e^u+e^{-u})^2} \\
&=\left(\frac{2}{e^u+e^{-u}}\right)^2 \\
&=\operatorname{sech}^2u
\end{aligned}
\]

Hence
\[
1-\tanh^2u=\operatorname{sech}^2u
\]

\item Let
\[
u=\operatorname{artanh}y
\]
so that
\[
y=\tanh u
\]

Using the exponential form of $\tanh u$,
\[
y=\frac{e^u-e^{-u}}{e^u+e^{-u}}
\]

Multiply top and bottom by $e^u$:
\[
y=\frac{e^{2u}-1}{e^{2u}+1}
\]

Now solve for $u$:
\[
\begin{aligned}
y(e^{2u}+1)&=e^{2u}-1 \\
ye^{2u}+y&=e^{2u}-1 \\
e^{2u}-ye^{2u}&=1+y \\
e^{2u}(1-y)&=1+y \\
e^{2u}&=\frac{1+y}{1-y}
\end{aligned}
\]

Since $|y|<1$, both $1+y$ and $1-y$ are positive, so taking logarithms is valid:
\[
2u=\ln\left(\frac{1+y}{1-y}\right)
\]

Therefore
\[
u=\frac12\ln\left(\frac{1+y}{1-y}\right)
\]

But $u=\operatorname{artanh}y$, so
\[
\operatorname{artanh}y=\frac12\ln\left(\frac{1+y}{1-y}\right)
\]

\item Let
\[
x=4\tanh u
\]

Then
\[
\frac{\mathrm{d}x}{\mathrm{d}u}=4\operatorname{sech}^2u
\qquad\Rightarrow\qquad
\mathrm{d}x=4\operatorname{sech}^2u\,\mathrm{d}u
\]

Also,
\[
\begin{aligned}
16-x^2
&=16-16\tanh^2u \\
&=16(1-\tanh^2u) \\
&=16\operatorname{sech}^2u
\end{aligned}
\]
using part (i).

So
\[
\begin{aligned}
\int \frac{1}{16-x^2}\,\mathrm{d}x
&=\int \frac{1}{16\operatorname{sech}^2u}\left(4\operatorname{sech}^2u\,\mathrm{d}u\right) \\
&=\frac14\int 1\,\mathrm{d}u \\
&=\frac{u}{4}+c
\end{aligned}
\]

Now $x=4\tanh u$ gives
\[
\tanh u=\frac{x}{4}
\]
so
\[
u=\operatorname{artanh}\frac{x}{4}
\]

Hence
\[
\int \frac{1}{16-x^2}\,\mathrm{d}x=\frac14\operatorname{artanh}\frac{x}{4}+c
\]
for $|x|<4$.

Now use part (ii) with $y=\dfrac{x}{4}$. Since $|x|<4$, we have $\left|\dfrac{x}{4}\right|<1$, so
\[
\operatorname{artanh}\frac{x}{4}
=\frac12\ln\left(\frac{1+\frac{x}{4}}{1-\frac{x}{4}}\right)
=\frac12\ln\left(\frac{4+x}{4-x}\right)
\]

Therefore
\[
\int \frac{1}{16-x^2}\,\mathrm{d}x
=\frac14\cdot \frac12\ln\left(\frac{4+x}{4-x}\right)+c
=\frac18\ln\left(\frac{4+x}{4-x}\right)+c
\]

So
\[
\int \frac{1}{16-x^2}\,\mathrm{d}x=\frac18\ln\left(\frac{4+x}{4-x}\right)+c
\]

\item First complete the square:
\[
\begin{aligned}
-t^2+6t+7
&=-(t^2-6t+9)+16 \\
&=16-(t-3)^2
\end{aligned}
\]

So
\[
\int_1^5 \frac{1}{-t^2+6t+7}\,\mathrm{d}t
=
\int_1^5 \frac{1}{16-(t-3)^2}\,\mathrm{d}t
\]

Let
\[
x=t-3
\]
Then $\mathrm{d}x=\mathrm{d}t$, and the limits become
\[
t=1 \Rightarrow x=-2,
\qquad
t=5 \Rightarrow x=2
\]

Hence
\[
\int_1^5 \frac{1}{-t^2+6t+7}\,\mathrm{d}t
=
\int_{-2}^{2}\frac{1}{16-x^2}\,\mathrm{d}x
\]

Using part (iii),
\[
\int \frac{1}{16-x^2}\,\mathrm{d}x
=
\frac18\ln\left(\frac{4+x}{4-x}\right)+c
\]

Therefore
\[
\begin{aligned}
\int_{-2}^{2}\frac{1}{16-x^2}\,\mathrm{d}x
&=
\left[\frac18\ln\left(\frac{4+x}{4-x}\right)\right]_{-2}^{2} \\
&=\frac18\ln\left(\frac{6}{2}\right)-\frac18\ln\left(\frac{2}{6}\right) \\
&=\frac18\ln 3-\frac18\ln\left(\frac13\right) \\
&=\frac18\ln 3+\frac18\ln 3 \\
&=\frac14\ln 3
\end{aligned}
\]

Hence
\[
\int_{1}^{5} \frac{1}{-t^2 + 6t + 7}\,\mathrm{d}t = \frac{1}{4}\ln 3
\]
\end{enumerate}
\end{tcolorbox}


\newpage

% ---- Question 4 ----
\questionitem
% [Source: _QBT___Solns__Integration_with_Inverse_Trig, Q4]

\begin{enumerate}[label=(\roman*), start=1]
  \item For $x>5$, find
  \[
  \int \frac{1}{\sqrt{x^2-6x+5}}\,\mathrm{d}x
  \] \\
  expressing your answer first in terms of $\operatorname{arcosh}$ and hence in logarithmic form. \marks{3} \\
  \item Use integration by substitution with $x-3=2\cosh u$ to show that, for $x>5$,
  \[
  \int \sqrt{x^2-6x+5}\,\mathrm{d}x=2\left(\frac{x-3}{2}\sqrt{\frac{(x-3)^2}{4}-1}-\ln\left|\frac{x-3}{2}+\sqrt{\frac{(x-3)^2}{4}-1}\right|\right)+c
  \]
  where $c$ is an arbitrary constant. \marks{7} \\
\end{enumerate}

\vspace*{10pt}

\begin{tcolorbox}[colback=white, colframe=scarlet, title={\textbf{Solution}}, breakable, grow to left by=6mm, grow to right by=10mm]
\begin{enumerate}[label=\textbf{(\roman*)}, start=1]
  \item First complete the square:
  \[
  x^2-6x+5=(x-3)^2-4
  \]

  Since $x>5$, we have $\dfrac{x-3}{2}>1$, so we may let
  \[
  x-3=2\cosh u \qquad (u>0)
  \]
  Then
  \[
  \frac{\mathrm{d}x}{\mathrm{d}u}=2\sinh u
  \qquad\Rightarrow\qquad
  \mathrm{d}x=2\sinh u\,\mathrm{d}u
  \]
  and
  \begin{align*}
  \sqrt{x^2-6x+5}
  &=\sqrt{(x-3)^2-4} \\
  &=\sqrt{4\cosh^2u-4} \\
  &=2\sqrt{\cosh^2u-1} \\
  &=2\sinh u
  \end{align*}
  since $u>0$.

  So
  \begin{align*}
  \int \frac{1}{\sqrt{x^2-6x+5}}\,\mathrm{d}x
  &=\int \frac{1}{2\sinh u}\,(2\sinh u\,\mathrm{d}u) \\
  &=\int 1\,\mathrm{d}u \\
  &=u+c
  \end{align*}

  Now
  \[
  \cosh u=\frac{x-3}{2}
  \]
  so
  \[
  u=\operatorname{arcosh}\left(\frac{x-3}{2}\right)
  \]

  Hence
  \[
  \int \frac{1}{\sqrt{x^2-6x+5}}\,\mathrm{d}x
  =\operatorname{arcosh}\left(\frac{x-3}{2}\right)+c
  \]

  Using the standard result
  \[
  \operatorname{arcosh} z=\ln\left(z+\sqrt{z^2-1}\right)
  \]
  we get
  \[
  \int \frac{1}{\sqrt{x^2-6x+5}}\,\mathrm{d}x
  =\ln\left(\frac{x-3}{2}+\sqrt{\frac{(x-3)^2}{4}-1}\right)+c
  \]

  So the answer is
  \[
  \operatorname{arcosh}\left(\frac{x-3}{2}\right)+c
  \]
  and hence
  \[
  \ln\left(\frac{x-3}{2}+\sqrt{\frac{(x-3)^2}{4}-1}\right)+c
  \]

  \item Let
  \[
  I=\int \sqrt{x^2-6x+5}\,\mathrm{d}x
  \]
  and use the substitution
  \[
  x-3=2\cosh u \qquad (x>5)
  \]

  Then, as above,
  \[
  \sqrt{x^2-6x+5}=2\sinh u
  \qquad\text{and}\qquad
  \mathrm{d}x=2\sinh u\,\mathrm{d}u
  \]

  So
  \begin{align*}
  I
  &=\int (2\sinh u)(2\sinh u)\,\mathrm{d}u \\
  &=4\int \sinh^2u\,\mathrm{d}u
  \end{align*}

  Use the identity
  \[
  \cosh 2u=2\sinh^2u+1
  \]
  so that
  \[
  2\sinh^2u=\cosh 2u-1
  \]

  Therefore
  \begin{align*}
  I
  &=2\int (\cosh 2u-1)\,\mathrm{d}u \\
  &=2\left(\frac{1}{2}\sinh 2u-u\right)+c \\
  &=\sinh 2u-2u+c
  \end{align*}

  Now express this in terms of $x$.

  Since
  \[
  \cosh u=\frac{x-3}{2}
  \]
  we have
  \[
  \sinh u=\sqrt{\cosh^2u-1}
  =\sqrt{\frac{(x-3)^2}{4}-1}
  \]
  because $u>0$.

  Also
  \begin{align*}
  \sinh 2u
  &=2\sinh u\cosh u \\
  &=2\left(\sqrt{\frac{(x-3)^2}{4}-1}\right)\left(\frac{x-3}{2}\right) \\
  &=(x-3)\sqrt{\frac{(x-3)^2}{4}-1}
  \end{align*}

  Hence
  \begin{align*}
  I
  &=(x-3)\sqrt{\frac{(x-3)^2}{4}-1}-2u+c \\
  &=2\left(\frac{x-3}{2}\sqrt{\frac{(x-3)^2}{4}-1}-u\right)+c
  \end{align*}

  From part (ii),
  \[
  u=\operatorname{arcosh}\left(\frac{x-3}{2}\right)
  =\ln\left|\frac{x-3}{2}+\sqrt{\frac{(x-3)^2}{4}-1}\right|
  \]

  Therefore
  \[
  \int \sqrt{x^2-6x+5}\,\mathrm{d}x
  =2\left(\frac{x-3}{2}\sqrt{\frac{(x-3)^2}{4}-1}-\ln\left|\frac{x-3}{2}+\sqrt{\frac{(x-3)^2}{4}-1}\right|\right)+c
  \]

  This is the required result.
\end{enumerate}
\end{tcolorbox}

\newpage
% ---- Question 5 ----
\questionitem
% [Source: _QBT___Solns__Integration_with_Inverse_Trig, Q5]

\[
g(x)=\frac{2x^2-3x+31}{x^3-3x^2+9x+13}
\]
\\
\begin{enumerate}[label=\textbf{(\alph*)}]
  \item Express $g(x)$ in the form \[\frac{P}{x+1}+\frac{Q}{x^2-4x+13}\] where $P$ and $Q$ are constants to be found. \marks{3} \\
  \item Find $\int g(x)\,\mathrm{d}x$, giving your answer in the form \[A\ln|x+1|+B\arctan\left(\frac{x-2}{3}\right)+c\] where $A$ and $B$ are constants to be found. \marks{4} \\
  \item Hence show that $\int_0^\infty g(x)\,\mathrm{d}x$ diverges. \marks{2} \\
\end{enumerate}

\vspace*{10pt}

\begin{tcolorbox}[colback=white, colframe=scarlet, title={\textbf{Solution}}, breakable, grow to left by=6mm, grow to right by=10mm]
\begin{enumerate}[label=\textbf{(\alph*)}]
  \item First factorise the denominator:
\[
x^3-3x^2+9x+13=(x+1)(x^2-4x+13)
\]

So
\[
g(x)=\frac{2x^2-3x+31}{(x+1)(x^2-4x+13)}
\]
and we write
\[
\frac{2x^2-3x+31}{(x+1)(x^2-4x+13)}
=\frac{P}{x+1}+\frac{Q}{x^2-4x+13}
\]

Multiplying through by $(x+1)(x^2-4x+13)$ gives
\[
2x^2-3x+31=P(x^2-4x+13)+Q(x+1)
\]

Expand the right-hand side:
\[
2x^2-3x+31=Px^2-4Px+13P+Qx+Q
\]
\[
=Px^2+(-4P+Q)x+(13P+Q)
\]

Now compare coefficients with
\[
2x^2-3x+31
\]

So
\[
P=2
\]
and
\[
-4P+Q=-3
\]

Substituting $P=2$:
\[
-8+Q=-3
\]
\[
Q=5
\]

Checking the constant term:
\[
13P+Q=13(2)+5=31
\]
which is correct.

Hence
\[
P=2,\qquad Q=5
\]

  \item From part (a),
\[
g(x)=\frac{2}{x+1}+\frac{5}{x^2-4x+13}
\]

Complete the square in the quadratic:
\[
x^2-4x+13=(x-2)^2+9
\]

So
\[
g(x)=\frac{2}{x+1}+\frac{5}{(x-2)^2+9}
\]

Therefore
\[
\int g(x)\,\mathrm{d}x=\int \frac{2}{x+1}\,\mathrm{d}x+\int \frac{5}{(x-2)^2+9}\,\mathrm{d}x
\]

For the first integral,
\[
\int \frac{2}{x+1}\,\mathrm{d}x=2\ln|x+1|
\]

For the second integral, use the standard result
\[
\int \frac{1}{u^2+a^2}\,\mathrm{d}u=\frac{1}{a}\arctan\!\left(\frac{u}{a}\right)+c
\]

Here $u=x-2$ and $a=3$, so
\[
\int \frac{5}{(x-2)^2+9}\,\mathrm{d}x
=5\int \frac{1}{(x-2)^2+3^2}\,\mathrm{d}x
=\frac{5}{3}\arctan\!\left(\frac{x-2}{3}\right)
\]

Hence
\[
\int g(x)\,\mathrm{d}x
=2\ln|x+1|+\frac{5}{3}\arctan\!\left(\frac{x-2}{3}\right)+c
\]

So
\[
A=2,\qquad B=\frac{5}{3}
\]

  \item To consider the improper integral, write
\[
\int_0^\infty g(x)\,\mathrm{d}x=\lim_{t\to\infty}\int_0^t g(x)\,\mathrm{d}x
\]

Using the antiderivative from part (b),
\begin{align*}
\int_0^t g(x)\,\mathrm{d}x
&=\left[2\ln|x+1|+\frac{5}{3}\arctan\!\left(\frac{x-2}{3}\right)\right]_0^t \\
&=2\ln(t+1)+\frac{5}{3}\arctan\!\left(\frac{t-2}{3}\right)
-\left(2\ln 1+\frac{5}{3}\arctan\!\left(-\frac23\right)\right) \\
&=2\ln(t+1)+\frac{5}{3}\arctan\!\left(\frac{t-2}{3}\right)+\frac{5}{3}\arctan\!\left(\frac23\right)
\end{align*}

Now as $t\to\infty$,
\[
\arctan\!\left(\frac{t-2}{3}\right)\to \frac{\pi}{2}
\]
so that term stays finite, but
\[
2\ln(t+1)\to\infty
\]

Therefore
\[
\int_0^\infty g(x)\,\mathrm{d}x
\]
diverges, in fact it diverges to $+\infty$.
\end{enumerate}
\end{tcolorbox}


\newpage

% ---- Question 6 ----
\questionitem
% [Source: _QBT___Solns__Integration_with_Inverse_Trig, Q6]

Show that
\[
\int_{1}^{4} \frac{1}{\sqrt{10 + 8x - 2x^2}}\,\mathrm{d}x = \frac{1}{\sqrt{2}}\left(\arcsin \left(\frac{2}{3}\right) + \arcsin \left(\frac{1}{3}\right)\right)
\]
\marks[-20pt]{6}

\vspace*{10pt}

\begin{tcolorbox}[colback=white, colframe=scarlet, title={\textbf{Solution}}, breakable, grow to left by=6mm, grow to right by=10mm]
Let
\[
I=\int_{1}^{4}\frac{1}{\sqrt{10+8x-2x^2}}\,\mathrm{d}x
\]

First complete the square in the expression under the root:
\begin{align*}
10+8x-2x^2
&=-2(x^2-4x-5)\\
&=-2\bigl((x-2)^2-9\bigr)\\
&=2\bigl(9-(x-2)^2\bigr)
\end{align*}

So
\[
I=\int_{1}^{4}\frac{1}{\sqrt{2\bigl(9-(x-2)^2\bigr)}}\,\mathrm{d}x
=\frac{1}{\sqrt{2}}\int_{1}^{4}\frac{1}{\sqrt{9-(x-2)^2}}\,\mathrm{d}x
\]

Now reduce this to the standard form
\[
\int \frac{1}{\sqrt{1-u^2}}\,\mathrm{d}u=\arcsin (u)+c
\]

Let
\[
u=\frac{x-2}{3}
\]
Then
\[
x-2=3u,\qquad \mathrm{d}x=3\,\mathrm{d}u
\]

Also, when $x=1$,
\[
u=\frac{1-2}{3}=-\frac13
\]
and when $x=4$,
\[
u=\frac{4-2}{3}=\frac23
\]

Substituting into the integral:
\begin{align*}
I
&=\frac{1}{\sqrt{2}}\int_{1}^{4}\frac{1}{\sqrt{9-(x-2)^2}}\,\mathrm{d}x\\
&=\frac{1}{\sqrt{2}}\int_{-1/3}^{2/3}\frac{3}{\sqrt{9-(3u)^2}}\,\mathrm{d}u\\
&=\frac{1}{\sqrt{2}}\int_{-1/3}^{2/3}\frac{3}{\sqrt{9-9u^2}}\,\mathrm{d}u\\
&=\frac{1}{\sqrt{2}}\int_{-1/3}^{2/3}\frac{3}{3\sqrt{1-u^2}}\,\mathrm{d}u\\
&=\frac{1}{\sqrt{2}}\int_{-1/3}^{2/3}\frac{1}{\sqrt{1-u^2}}\,\mathrm{d}u
\end{align*}

Using the standard result:
\begin{align*}
I
&=\frac{1}{\sqrt{2}}\left[\arcsin (u)\right]_{-1/3}^{2/3}\\
&=\frac{1}{\sqrt{2}}\left(\arcsin \left(\frac23\right)-\arcsin \left(-\frac13\right)\right)
\end{align*}

Since $\arcsin (-a)=-\arcsin (a)$,
\[
I=\frac{1}{\sqrt{2}}\left(\arcsin \left(\frac23\right)+\arcsin \left(\frac13\right)\right)
\]

Hence
\[
\int_{1}^{4} \frac{1}{\sqrt{10 + 8x - 2x^2}}\,\mathrm{d}x
= \frac{1}{\sqrt{2}}\left(\arcsin \left(\frac{2}{3}\right) + \arcsin \left(\frac{1}{3}\right)\right)
\]
\end{tcolorbox}


\newpage

% ---- Question 7 ----
\questionitem
% [Source: _QBT___Solns__Integration_with_Inverse_Trig, Q7]

\begin{enumerate}[label=\textbf{(\alph*)}]
  \item Explain why
  \[
  \int_{0}^{\infty} \frac{1}{x^2 - 6x + 18} \, \mathrm{d}x
  \]
  is an improper integral. \marks{1} \\
  \item Show that
  \[
  \int_{0}^{\infty} \frac{1}{x^2 - 6x + 18} \, \mathrm{d}x = k\pi
  \]
  where $k$ is a constant to be determined. \marks{4} \\
\end{enumerate}

\vspace*{10pt}

\begin{tcolorbox}[colback=white, colframe=scarlet, title={\textbf{Solution}}, breakable, grow to left by=6mm, grow to right by=10mm]
\begin{enumerate}[label=\textbf{(\alph*)}]
  \item The integral is improper because the interval of integration is unbounded: the upper limit is $\infty$.

  So it is defined as
  \[
  \int_{0}^{\infty} \frac{1}{x^2-6x+18}\,\mathrm{d}x
  =
  \lim_{b\to\infty}\int_{0}^{b}\frac{1}{x^2-6x+18}\,\mathrm{d}x
  \]

  \item Let
  \[
  I=\int_{0}^{\infty}\frac{1}{x^2-6x+18}\,\mathrm{d}x
  \]

  First complete the square in the denominator:
  \[
  x^2-6x+18=(x-3)^2+9=(x-3)^2+3^2
  \]

  Hence
  \[
  I=\int_{0}^{\infty}\frac{1}{(x-3)^2+3^2}\,\mathrm{d}x
  \]

  Since this is an improper integral, write it as a limit:
  \[
  I=\lim_{b\to\infty}\int_{0}^{b}\frac{1}{(x-3)^2+3^2}\,\mathrm{d}x
  \]

  Using the standard result
  \[
  \int \frac{1}{u^2+a^2}\,\mathrm{d}u=\frac{1}{a}\arctan\!\left(\frac{u}{a}\right)+c
  \]
  with $u=x-3$ and $a=3$, we get
  \[
  \int \frac{1}{(x-3)^2+3^2}\,\mathrm{d}x
  =
  \frac{1}{3}\arctan\!\left(\frac{x-3}{3}\right)+c
  \]

  Therefore
  \begin{align*}
  I
  &=\lim_{b\to\infty}\left[\frac{1}{3}\arctan\!\left(\frac{x-3}{3}\right)\right]_{0}^{b} \\
  &=\frac{1}{3}\lim_{b\to\infty}\left(\arctan\!\left(\frac{b-3}{3}\right)-\arctan(-1)\right)
  \end{align*}

  Now
  \[
  \arctan\!\left(\frac{b-3}{3}\right)\to \frac{\pi}{2}
  \quad\text{as } b\to\infty
  \]
  and
  \[
  \arctan(-1)=-\frac{\pi}{4}
  \]

  So
  \begin{align*}
  I
  &=\frac{1}{3}\left(\frac{\pi}{2}-\left(-\frac{\pi}{4}\right)\right) \\
  &=\frac{1}{3}\left(\frac{3\pi}{4}\right) \\
  &=\frac{\pi}{4}
  \end{align*}

  Hence
  \[
  \int_{0}^{\infty}\frac{1}{x^2-6x+18}\,\mathrm{d}x=\frac{\pi}{4}=k\pi
  \]
  so
  \[
  k=\frac{1}{4}
  \]
\end{enumerate}
\end{tcolorbox}


\newpage

% ---- Question 8 ----
\questionitem
% [Source: _QBT___Solns__Integration_with_Inverse_Trig, Q8]

Evaluate
\[
\int_{0}^{4} \frac{1}{x^2-4x+8}\, \mathrm{d}x
\]
giving your answer as an exact multiple of $\pi$. \marks{8}

\vspace*{10pt}

\begin{tcolorbox}[colback=white, colframe=scarlet, title={\textbf{Solution}}, breakable, grow to left by=6mm, grow to right by=10mm]
First complete the square in the denominator:
\[
x^2-4x+8=(x^2-4x+4)+4=(x-2)^2+4=(x-2)^2+2^2
\]

So the integral becomes
\[
\int_0^4 \frac{1}{x^2-4x+8}\,\mathrm{d}x
=
\int_0^4 \frac{1}{(x-2)^2+2^2}\,\mathrm{d}x
\]

Now use the standard result from the formula booklet:
\[
\int \frac{1}{u^2+a^2}\,\mathrm{d}u=\frac{1}{a}\arctan\!\left(\frac{u}{a}\right)+c
\]

Here we can take
\[
u=x-2 \qquad \text{so} \qquad \mathrm{d}u=\mathrm{d}x
\]

Therefore
\[
\int_0^4 \frac{1}{(x-2)^2+2^2}\,\mathrm{d}x
=
\left[\frac{1}{2}\arctan\!\left(\frac{x-2}{2}\right)\right]_0^4
\]

Now substitute the limits:
\begin{align*}
\left[\frac{1}{2}\arctan\!\left(\frac{x-2}{2}\right)\right]_0^4
&=
\frac{1}{2}\arctan\!\left(\frac{4-2}{2}\right)
-
\frac{1}{2}\arctan\!\left(\frac{0-2}{2}\right) \\
&=
\frac{1}{2}\arctan(1)-\frac{1}{2}\arctan(-1) \\
&=
\frac{1}{2}\cdot\frac{\pi}{4}-\frac{1}{2}\cdot\left(-\frac{\pi}{4}\right) \\
&=
\frac{1}{2}\left(\frac{\pi}{4}+\frac{\pi}{4}\right) \\
&=
\frac{1}{2}\cdot\frac{\pi}{2} \\
&=
\frac{\pi}{4}
\end{align*}

Hence
\[
\int_{0}^{4} \frac{1}{x^2-4x+8}\,\mathrm{d}x=\frac{\pi}{4}
\]

So the exact value, as a multiple of $\pi$, is $\frac{\pi}{4}$.
\end{tcolorbox}


\newpage

% ---- Question 9 ----
\questionitem
% [Source: _QBT___Solns__Integration_with_Inverse_Trig, Q9]

\[f(x)=\frac{6x+1}{4x^2-8x+13}\] \\

\begin{enumerate}[label=\textbf{(\alph*)}]
  \item Find $\int f(x)\,\mathrm{d}x$, giving your answer in the form \[A\ln(4x^2-8x+13)+B\arctan\left(\frac{2x-2}{3}\right)+c\] where $c$ is an arbitrary constant and $A$ and $B$ are constants to be found. \marks{4} \\
  \item Hence show that the mean value of $f(x)$ over the interval $\left[1,\frac{5}{2}\right]$ is \[\frac{1}{2}\ln 2+\frac{7\pi}{36}\] \marks[-20pt]{3} 
  \item Hence write down the mean value of $3-2f(x)$ over the interval $\left[1,\frac{5}{2}\right]$. \marks{1} \\
\end{enumerate}

\vspace*{10pt}

\begin{tcolorbox}[colback=white, colframe=scarlet, title={\textbf{Solution}}, breakable, grow to left by=6mm, grow to right by=10mm]
\begin{enumerate}[label=\textbf{(\alph*)}]
  \item Let
\[
I=\int \frac{6x+1}{4x^2-8x+13}\,\mathrm{d}x
\]

Write the numerator in terms of the derivative of the denominator:
\[
6x+1=\frac34(8x-8)+7
\]

So
\begin{align*}
I
&=\frac34\int \frac{8x-8}{4x^2-8x+13}\,\mathrm{d}x
+7\int \frac{1}{4x^2-8x+13}\,\mathrm{d}x
\end{align*}

For the first integral,
\[
\int \frac{g'(x)}{g(x)}\,\mathrm{d}x=\ln|g(x)|+c
\]
and here
\[
4x^2-8x+13=(2x-2)^2+9>0
\]
so no modulus is needed. Hence
\[
\frac34\int \frac{8x-8}{4x^2-8x+13}\,\mathrm{d}x
=\frac34\ln(4x^2-8x+13)
\]

For the second integral, complete the square:
\[
4x^2-8x+13=(2x-2)^2+9
\]

Let \(u=2x-2\), so \(\mathrm{d}u=2\,\mathrm{d}x\) and \(\mathrm{d}x=\frac12\,\mathrm{d}u\). Then
\begin{align*}
7\int \frac{1}{(2x-2)^2+9}\,\mathrm{d}x
&=7\int \frac{1}{u^2+9}\cdot \frac12\,\mathrm{d}u \\
&=\frac72 \int \frac{1}{u^2+9}\,\mathrm{d}u
\end{align*}

Using the standard result
\[
\int \frac{1}{a^2+u^2}\,\mathrm{d}u=\frac1a\arctan\left(\frac{u}{a}\right)+c
\]
with \(a=3\),
\begin{align*}
\frac72 \int \frac{1}{u^2+9}\,\mathrm{d}u
&=\frac72 \cdot \frac13 \arctan\left(\frac{u}{3}\right) \\
&=\frac76 \arctan\left(\frac{u}{3}\right)
\end{align*}

Substituting back \(u=2x-2\),
\[
\int f(x)\,\mathrm{d}x
=\frac34\ln(4x^2-8x+13)+\frac76\arctan\left(\frac{2x-2}{3}\right)+c
\]

Therefore
\[
A=\frac34,\qquad B=\frac76
\]

  \item From part (a), an antiderivative is
\[
F(x)=\frac34\ln(4x^2-8x+13)+\frac76\arctan\left(\frac{2x-2}{3}\right)
\]

The mean value of \(f(x)\) over \(\left[1,\frac52\right]\) is
\[
\frac{1}{\frac52-1}\int_1^{5/2} f(x)\,\mathrm{d}x
=\frac23\bigl(F(5/2)-F(1)\bigr)
\]

Now
\begin{align*}
F\left(\frac52\right)
&=\frac34\ln\left(4\cdot\frac{25}{4}-8\cdot\frac52+13\right)
+\frac76\arctan\left(\frac{2\cdot\frac52-2}{3}\right) \\
&=\frac34\ln(18)+\frac76\arctan(1) \\
&=\frac34\ln(18)+\frac{7\pi}{24}
\end{align*}

Also
\begin{align*}
F(1)
&=\frac34\ln(4-8+13)+\frac76\arctan(0) \\
&=\frac34\ln(9)
\end{align*}

So
\begin{align*}
\frac23\bigl(F(5/2)-F(1)\bigr)
&=\frac23\left(\frac34\ln\frac{18}{9}+\frac{7\pi}{24}\right) \\
&=\frac23\left(\frac34\ln 2+\frac{7\pi}{24}\right) \\
&=\frac12\ln 2+\frac{7\pi}{36}
\end{align*}

Hence the mean value is
\[
\frac12\ln 2+\frac{7\pi}{36}
\]

  \item Using linearity of the mean value,
\[
\text{mean value of }(3-2f(x))
=3-2\times \text{mean value of }f(x)
\]

So
\begin{align*}
\text{mean value of }(3-2f(x))
&=3-2\left(\frac12\ln 2+\frac{7\pi}{36}\right) \\
&=3-\ln 2-\frac{7\pi}{18}
\end{align*}

Therefore the mean value is
\[
3-\ln 2-\frac{7\pi}{18}
\]
\end{enumerate}
\end{tcolorbox}


\newpage

% ---- Question 10 ----
\questionitem
% [Source: _QBT___Solns__Integration_with_Inverse_Trig, Q10]

\[
g(x)=\frac{4x-1}{x^2-4x+8}
\] \\
Show that
\[
\int g(x)\,\mathrm{d}x=A\arctan\left(\frac{x-2}{2}\right)+B\ln(x^2-4x+8)+c
\]
where $c$ is an arbitrary constant and $A$ and $B$ are constants to be found. \marks{5} \\

\vspace*{10pt}

\begin{tcolorbox}[colback=white, colframe=scarlet, title={\textbf{Solution}}, breakable, grow to left by=6mm, grow to right by=10mm]
First complete the square in the denominator:
\[
x^2-4x+8=(x-2)^2+4
\]

Now rewrite the numerator so that part of it is the derivative of the denominator:
\[
4x-1=2(2x-4)+7
\]

So
\begin{align*}
g(x)
&=\frac{4x-1}{x^2-4x+8} \\
&=\frac{2(2x-4)+7}{x^2-4x+8} \\
&=2\frac{2x-4}{x^2-4x+8}+\frac{7}{(x-2)^2+4}
\end{align*}

Hence
\begin{align*}
\int g(x)\,\mathrm{d}x
&=2\int \frac{2x-4}{x^2-4x+8}\,\mathrm{d}x
+7\int \frac{1}{(x-2)^2+4}\,\mathrm{d}x
\end{align*}

For the first integral, let
\[
u=x^2-4x+8
\qquad\text{so}\qquad
\frac{\mathrm{d}u}{\mathrm{d}x}=2x-4
\]
Then
\[
2\int \frac{2x-4}{x^2-4x+8}\,\mathrm{d}x
=2\int \frac{1}{u}\,\mathrm{d}u
=2\ln u
=2\ln(x^2-4x+8)
\]

For the second integral, use the standard result
\[
\int \frac{1}{a^2+u^2}\,\mathrm{d}u
=\frac{1}{a}\arctan\left(\frac{u}{a}\right)+c
\]
Here \(u=x-2\) and \(a=2\), so
\[
7\int \frac{1}{(x-2)^2+4}\,\mathrm{d}x
=7\cdot \frac{1}{2}\arctan\left(\frac{x-2}{2}\right)
\]

Therefore
\begin{align*}
\int g(x)\,\mathrm{d}x
&=2\ln(x^2-4x+8)+\frac{7}{2}\arctan\left(\frac{x-2}{2}\right)+c
\end{align*}

This is of the form
\[
\int g(x)\,\mathrm{d}x
=A\arctan\left(\frac{x-2}{2}\right)+B\ln(x^2-4x+8)+c
\]
with
\[
A=\frac{7}{2},
\qquad
B=2
\]

So the required constants are \(A=\frac{7}{2}\) and \(B=2\).
\end{tcolorbox}


\newpage

% ---- Question 11 ----
\questionitem
% [Source: _QBT___Solns__Integration_with_Inverse_Trig, Q11]

You are given that $y=\operatorname{arcosh}x$, where $x\ge 1$. \\

\begin{enumerate}[label=(\roman*), start=1]
  \item Show that
  \[
  \frac{\mathrm{d}y}{\mathrm{d}x}=\frac{1}{\sqrt{x^2-1}}
  \]
  \marks[-20pt]{4}
  \item Show, by integrating the result in part (i), that
  \[
  y=\ln\left(x+\sqrt{x^2-1}\right)
  \]
  \marks[-20pt]{4}
  \item Hence show that
  \[
  \int_{15/8}^{5/2} \frac{1}{\sqrt{4x^2-9}}\,\mathrm{d}x=\frac{1}{2}\left(\operatorname{arcosh}\frac{5}{3}-\operatorname{arcosh}\frac{5}{4}\right)
  \]
  Express this answer in logarithmic form. \marks{4} \\
\end{enumerate}

\vspace*{10pt}

\begin{tcolorbox}[colback=white, colframe=scarlet, title={\textbf{Solution}}, breakable, grow to left by=6mm, grow to right by=10mm]
\begin{enumerate}[label=\textbf{(\roman*)}]
\item Since \(y=\operatorname{arcosh}x\), we have
\[
\cosh y=x
\]
Differentiate implicitly with respect to \(x\):
\[
\sinh y\frac{\mathrm{d}y}{\mathrm{d}x}=1
\]
So
\[
\frac{\mathrm{d}y}{\mathrm{d}x}=\frac{1}{\sinh y}
\]
Using the identity
\[
\cosh^2 y-\sinh^2 y=1
\]
gives
\[
\sinh^2 y=\cosh^2 y-1=x^2-1
\]
Also \(y=\operatorname{arcosh}x\) with \(x\ge 1\), so \(y\ge 0\), hence \(\sinh y\ge 0\). Therefore
\[
\sinh y=\sqrt{x^2-1}
\]
and so
\[
\frac{\mathrm{d}y}{\mathrm{d}x}=\frac{1}{\sqrt{x^2-1}}
\]
as required.

\item Integrating the result from part (ii),
\[
y=\int \frac{1}{\sqrt{x^2-1}}\,\mathrm{d}x
\]
Use the standard result
\[
\int \frac{1}{\sqrt{x^2-a^2}}\,\mathrm{d}x=\ln\left|x+\sqrt{x^2-a^2}\right|+C
\]
With \(a=1\),
\[
y=\ln\left|x+\sqrt{x^2-1}\right|+C
\]
Since \(x\ge 1\), the quantity \(x+\sqrt{x^2-1}\) is positive, so
\[
y=\ln\left(x+\sqrt{x^2-1}\right)+C
\]
Now when \(x=1\),
\[
y=\operatorname{arcosh}1=0
\]
Hence
\[
0=\ln\left(1+\sqrt{1^2-1}\right)+C=\ln 1+C=C
\]
So \(C=0\), and therefore
\[
y=\ln\left(x+\sqrt{x^2-1}\right)
\]

\medskip

\emph{Alternative method (without quoting the standard result):}

Let
\[
I=\int \frac{1}{\sqrt{x^2-1}}\,\mathrm{d}x
\]
Multiply the integrand by
\[
\frac{x+\sqrt{x^2-1}}{x+\sqrt{x^2-1}}
\]
to obtain
\begin{align*}
I
&=\int \frac{x+\sqrt{x^2-1}}{\sqrt{x^2-1}\left(x+\sqrt{x^2-1}\right)}\,\mathrm{d}x \\[4pt]
&=\int \frac{1+\dfrac{x}{\sqrt{x^2-1}}}{x+\sqrt{x^2-1}}\,\mathrm{d}x
\end{align*}
Now
\[
\frac{\mathrm{d}}{\mathrm{d}x}\!\left(x+\sqrt{x^2-1}\right)
=1+\frac{x}{\sqrt{x^2-1}}
\]
so
\[
I=\int \frac{f'(x)}{f(x)}\,\mathrm{d}x
\quad\text{with}\quad
f(x)=x+\sqrt{x^2-1}
\]
Hence
\[
I=\ln\left(x+\sqrt{x^2-1}\right)+C
\]
since \(x\ge 1\) implies \(x+\sqrt{x^2-1}>0\)

Therefore
\[
y=\ln\left(x+\sqrt{x^2-1}\right)+C
\]
Now when \(x=1\),
\[
y=\operatorname{arcosh}1=0
\]
so
\[
0=\ln\left(1+\sqrt{1^2-1}\right)+C=\ln 1+C=C
\]
Thus \(C=0\), giving
\[
y=\ln\left(x+\sqrt{x^2-1}\right)
\] \\

\item Let
\[
u=\frac{2x}{3}
\]
Then
\[
\mathrm{d}x=\frac{3}{2}\,\mathrm{d}u
\]
and
\[
\sqrt{4x^2-9}=\sqrt{9u^2-9}=3\sqrt{u^2-1}
\]
Also the limits become
\[
x=\frac{15}{8}\Rightarrow u=\frac{5}{4},
\qquad
x=\frac{5}{2}\Rightarrow u=\frac{5}{3}
\]
So
\begin{align*}
\int_{15/8}^{5/2}\frac{1}{\sqrt{4x^2-9}}\,\mathrm{d}x
&=\int_{5/4}^{5/3}\frac{1}{3\sqrt{u^2-1}}\cdot \frac{3}{2}\,\mathrm{d}u\\
&=\frac12\int_{5/4}^{5/3}\frac{1}{\sqrt{u^2-1}}\,\mathrm{d}u
\end{align*}
From part (ii), \(\dfrac{\mathrm{d}}{\mathrm{d}u}(\operatorname{arcosh}u)=\dfrac{1}{\sqrt{u^2-1}}\), so
\[
\int_{15/8}^{5/2}\frac{1}{\sqrt{4x^2-9}}\,\mathrm{d}x
=\frac12\left[\operatorname{arcosh}u\right]_{5/4}^{5/3}
=\frac12\left(\operatorname{arcosh}\frac53-\operatorname{arcosh}\frac54\right)
\]
which is the required result.

Now express this in logarithmic form using part (iii):
\[
\operatorname{arcosh}t=\ln\left(t+\sqrt{t^2-1}\right)
\]
Hence
\[
\operatorname{arcosh}\frac53
=\ln\left(\frac53+\sqrt{\frac{25}{9}-1}\right)
=\ln\left(\frac53+\frac43\right)
=\ln 3
\]
and
\[
\operatorname{arcosh}\frac54
=\ln\left(\frac54+\sqrt{\frac{25}{16}-1}\right)
=\ln\left(\frac54+\frac34\right)
=\ln 2
\]
Therefore
\[
\int_{15/8}^{5/2}\frac{1}{\sqrt{4x^2-9}}\,\mathrm{d}x
=\frac12(\ln 3-\ln 2)
=\frac12\ln\left(\frac32\right)
\]
So the logarithmic form is
\[
\frac12\ln\left(\frac32\right)
\]
\end{enumerate}
\end{tcolorbox}


\newpage

% ---- Question 12 ----
\questionitem
% [Source: _QBT___Solns__Integration_with_Inverse_Trig, Q12]

Determine each of the following integrals. \\
\begin{enumerate}[label=(\roman*)]
  \item
  \[
  \int \frac{1}{9x^2+12x+20}\,\mathrm{d}x
  \]
  \marks[-20pt]{4}
  \item
  \[
  \int \frac{1}{\sqrt{45+12x-9x^2}}\,\mathrm{d}x
  \]
  \marks[-20pt]{4}
\end{enumerate}

\vspace*{10pt}

\begin{tcolorbox}[colback=white, colframe=scarlet, title={\textbf{Solution}}, breakable, grow to left by=6mm, grow to right by=10mm]
\begin{enumerate}[label=\textbf{(\roman*)}]
\item Complete the square in the denominator:
\begin{align*}
9x^2+12x+20
&=9\left(x^2+\frac{4}{3}x\right)+20 \\
&=9\left[\left(x+\frac23\right)^2-\frac49\right]+20 \\
&=9\left(x+\frac23\right)^2-4+20 \\
&=9\left(x+\frac23\right)^2+16 \\
&=9\left[\left(x+\frac23\right)^2+\left(\frac43\right)^2\right]
\end{align*}

So
\begin{align*}
\int \frac{1}{9x^2+12x+20}\,\mathrm{d}x
&=\frac19\int \frac{1}{\left(x+\frac23\right)^2+\left(\frac43\right)^2}\,\mathrm{d}x
\end{align*}

Let
\[
u=x+\frac23
\]
so that $\mathrm{d}u=\mathrm{d}x$. Then
\begin{align*}
\int \frac{1}{9x^2+12x+20}\,\mathrm{d}x
&=\frac19\int \frac{1}{u^2+\left(\frac43\right)^2}\,\mathrm{d}u
\end{align*}

Using the standard result
\[
\int \frac{1}{x^2+a^2}\,\mathrm{d}x=\frac{1}{a}\arctan\!\left(\frac{x}{a}\right)+C
\]
with $a=\frac43$,
\begin{align*}
\int \frac{1}{9x^2+12x+20}\,\mathrm{d}x
&=\frac19 \cdot \frac{1}{4/3}\arctan\!\left(\frac{u}{4/3}\right)+C \\
&=\frac{1}{12}\arctan\!\left(\frac{3u}{4}\right)+C \\
&=\frac{1}{12}\arctan\!\left(\frac{3x+2}{4}\right)+C
\end{align*}

Hence,
\[
\int \frac{1}{9x^2 + 12x + 20}\,\mathrm{d}x=\frac{1}{12}\arctan\!\left(\frac{3x+2}{4}\right)+C
\]

\item Rewrite the expression under the root in completed-square form:
\begin{align*}
45+12x-9x^2
&=-9\left(x^2-\frac43x-5\right) \\
&=-9\left[\left(x-\frac23\right)^2-\frac49-5\right] \\
&=-9\left[\left(x-\frac23\right)^2-\frac{49}{9}\right] \\
&=49-9\left(x-\frac23\right)^2 \\
&=49-(3x-2)^2
\end{align*}

So
\begin{align*}
\int \frac{1}{\sqrt{45+12x-9x^2}}\,\mathrm{d}x
&=\int \frac{1}{\sqrt{49-(3x-2)^2}}\,\mathrm{d}x
\end{align*}

Let
\[
u=3x-2
\]
Then
\[
\frac{\mathrm{d}u}{\mathrm{d}x}=3
\qquad\text{so}\qquad
\mathrm{d}x=\frac13\,\mathrm{d}u
\]

Hence
\begin{align*}
\int \frac{1}{\sqrt{45+12x-9x^2}}\,\mathrm{d}x
&=\frac13\int \frac{1}{\sqrt{49-u^2}}\,\mathrm{d}u
\end{align*}

Using the standard result
\[
\int \frac{1}{\sqrt{a^2-x^2}}\,\mathrm{d}x=\arcsin\!\left(\frac{x}{a}\right)+C
\]
with $a=7$,
\begin{align*}
\int \frac{1}{\sqrt{45+12x-9x^2}}\,\mathrm{d}x
&=\frac13\arcsin\!\left(\frac{u}{7}\right)+C \\
&=\frac13\arcsin\!\left(\frac{3x-2}{7}\right)+C
\end{align*}

Hence,
\[
\int \frac{1}{\sqrt{45 + 12x - 9x^2}}\,\mathrm{d}x=\frac13\arcsin\!\left(\frac{3x-2}{7}\right)+C
\]
\end{enumerate}
\end{tcolorbox}


\newpage

% ---- Question 13 ----
\questionitem
% [Source: _QBT___Solns__Integration_with_Inverse_Trig, Q13]

Without using a calculator, find \\
\begin{enumerate}[label=\textbf{(\alph*)}]
  \item
  \[
  \int_{-1/5}^{1} \frac{1}{25x^2-20x+13}\,\mathrm{d}x
  \] \\
  giving your answer as a multiple of $\pi$ \marks{5} \\
  \item
  \[
  \int_{-1/5}^{1} \frac{1}{\sqrt{25x^2-20x+13}}\,\mathrm{d}x
  \] \\
  giving your answer in the form $p\ln(q+r\sqrt{2})$, where $p$, $q$ and $r$ are rational numbers to be found. \marks{7} \\
\end{enumerate}

\vspace*{10pt}

\begin{tcolorbox}[colback=white, colframe=scarlet, title={\textbf{Solution}}, breakable, grow to left by=6mm, grow to right by=10mm]
\begin{enumerate}[label=\textbf{(\alph*)}]
  \item Complete the square in the denominator:
  \begin{align*}
  25x^2-20x+13
  &= (5x-2)^2+9
  \end{align*}

  So
  \[
  \int_{-\frac15}^{1}\frac{1}{25x^2-20x+13}\,\mathrm{d}x
  = \int_{-\frac15}^{1}\frac{1}{(5x-2)^2+9}\,\mathrm{d}x
  \]

  Let
  \[
  u=\frac{5x-2}{3}
  \]
  Then
  \[
  5x-2=3u,
  \qquad \mathrm{d}x=\frac35\,\mathrm{d}u
  \]

  Change the limits:
  \[
  x=-\frac15 \;\Rightarrow\; u=\frac{-1-2}{3}=-1,
  \qquad
  x=1 \;\Rightarrow\; u=\frac{5-2}{3}=1
  \]

  Hence
  \begin{align*}
  \int_{-\frac15}^{1}\frac{1}{(5x-2)^2+9}\,\mathrm{d}x
  &= \int_{-1}^{1}\frac{1}{9u^2+9}\cdot \frac35\,\mathrm{d}u \\
  &= \frac{1}{15}\int_{-1}^{1}\frac{1}{1+u^2}\,\mathrm{d}u
  \end{align*}

  Using the standard result
  \[
  \int \frac{1}{1+u^2}\,\mathrm{d}u=\arctan u + c
  \]

  we get
  \begin{align*}
  \frac{1}{15}\int_{-1}^{1}\frac{1}{1+u^2}\,\mathrm{d}u
  &= \frac{1}{15}\left[\arctan u\right]_{-1}^{1} \\
  &= \frac{1}{15}\left(\arctan 1-\arctan(-1)\right) \\
  &= \frac{1}{15}\left(\frac{\pi}{4}-\left(-\frac{\pi}{4}\right)\right) \\
  &= \frac{1}{15}\cdot \frac{\pi}{2} \\
  &= \frac{\pi}{30}
  \end{align*}

  Therefore, the value of the integral is
  \[
  \frac{\pi}{30}
  \]

  \item As in part (a),
  \[
  25x^2-20x+13=(5x-2)^2+9
  \]

  So
  \[
  \int_{-\frac15}^{1}\frac{1}{\sqrt{25x^2-20x+13}}\,\mathrm{d}x
  = \int_{-\frac15}^{1}\frac{1}{\sqrt{(5x-2)^2+9}}\,\mathrm{d}x
  \]

  Let
  \[
  u=\frac{5x-2}{3}
  \]
  Then
  \[
  5x-2=3u,
  \qquad \mathrm{d}x=\frac35\,\mathrm{d}u
  \]
  and the limits again become
  \[
  x=-\frac15 \Rightarrow u=-1,
  \qquad
  x=1 \Rightarrow u=1
  \]

  Therefore
  \begin{align*}
  \int_{-\frac15}^{1}\frac{1}{\sqrt{(5x-2)^2+9}}\,\mathrm{d}x
  &= \int_{-1}^{1}\frac{1}{\sqrt{9u^2+9}}\cdot \frac35\,\mathrm{d}u \\
  &= \frac15\int_{-1}^{1}\frac{1}{\sqrt{u^2+1}}\,\mathrm{d}u
  \end{align*}

  Using the standard result
  \[
  \int \frac{1}{\sqrt{u^2+a^2}}\,\mathrm{d}u
  = \ln\left|u+\sqrt{u^2+a^2}\right|+c
  \]
  with $a=1$, we have
  \[
  \int \frac{1}{\sqrt{u^2+1}}\,\mathrm{d}u
  = \ln\left|u+\sqrt{u^2+1}\right|+c
  \]

  Hence
  \begin{align*}
  \frac15\int_{-1}^{1}\frac{1}{\sqrt{u^2+1}}\,\mathrm{d}u
  &= \frac15\left[\ln\left(u+\sqrt{u^2+1}\right)\right]_{-1}^{1} \\
  &= \frac15\left(\ln(1+\sqrt2)-\ln(\sqrt2-1)\right) \\
  &= \frac15\ln\left(\frac{1+\sqrt2}{\sqrt2-1}\right)
  \end{align*}

  Now simplify the fraction by rationalising:
  \begin{align*}
  \frac{1+\sqrt2}{\sqrt2-1}
  &= \frac{(1+\sqrt2)(\sqrt2+1)}{(\sqrt2-1)(\sqrt2+1)} \\
  &= \frac{(1+\sqrt2)^2}{2-1} \\
  &= (1+\sqrt2)^2 \\
  &= 1+2\sqrt2+2 \\
  &= 3+2\sqrt2
  \end{align*}

  So the integral is
  \[
  \frac15\ln(3+2\sqrt2)
  \]

  Therefore, in the form $p\ln(q+r\sqrt2)$,
  \[
  p=\frac15,\qquad q=3,\qquad r=2
  \]
\end{enumerate}
\end{tcolorbox}


\newpage

% ---- Question 14 ----
\questionitem
% [Source: _QBT___Solns__Integration_with_Inverse_Trig, Q14]

\begin{enumerate}[label=(\roman*)]
  \item Given that $\tanh y=x$, where $|x|<1$, show that
  \[
  y=\frac{1}{2}\ln\left(\frac{1+x}{1-x}\right)
  \]
  and hence that
  \[
  \operatorname{artanh}x=\frac{1}{2}\ln\left(\frac{1+x}{1-x}\right)
  \]
  \marks[-20pt]{6}
  \item Use your result in part (i), or otherwise, to find
  \[
  \int_{1/5}^{2/5} \frac{1}{4-9x^2}\,\mathrm{d}x
  \]
  expressing your answer in an exact logarithmic form. \marks{6} \\
\end{enumerate}

\vspace*{10pt}

\begin{tcolorbox}[colback=white, colframe=scarlet, title={\textbf{Solution}}, breakable, grow to left by=6mm, grow to right by=10mm]
\begin{enumerate}[label=\textbf{(\roman*)}]
\item Starting from the definition of the hyperbolic tangent,
\[
\tanh y=\frac{e^y-e^{-y}}{e^y+e^{-y}}
\]
Multiplying numerator and denominator by $e^y$ gives
\[
\tanh y=\frac{e^{2y}-1}{e^{2y}+1}
\]
Given that $\tanh y=x$, we have
\begin{align*}
x&=\frac{e^{2y}-1}{e^{2y}+1} \\
x(e^{2y}+1)&=e^{2y}-1 \\
xe^{2y}+x&=e^{2y}-1 \\
e^{2y}-xe^{2y}&=x+1 \\
e^{2y}(1-x)&=1+x
\end{align*}
Since $|x|<1$, we have $1-x>0$, so we may divide through by $1-x$:
\[
e^{2y}=\frac{1+x}{1-x}
\]
Now take natural logarithms:
\[
2y=\ln\left(\frac{1+x}{1-x}\right)
\]
Hence
\[
y=\frac{1}{2}\ln\left(\frac{1+x}{1-x}\right)
\]
So the required result is shown.

Since $\operatorname{artanh}x$ is the inverse of $\tanh y$, this means
\[
\operatorname{artanh}x=\frac{1}{2}\ln\left(\frac{1+x}{1-x}\right)
\]

\item From part (i),
\[
\operatorname{artanh}u=\frac12\ln\left(\frac{1+u}{1-u}\right)
\]
for $|u|<1$.

Differentiating,
\begin{align*}
\frac{\mathrm{d}}{\mathrm{d}u}\bigl(\operatorname{artanh}u\bigr)
&=\frac12\frac{\mathrm{d}}{\mathrm{d}u}\left[\ln(1+u)-\ln(1-u)\right] \\
&=\frac12\left(\frac{1}{1+u}+\frac{1}{1-u}\right) \\
&=\frac12\left(\frac{(1-u)+(1+u)}{1-u^2}\right) \\
&=\frac{1}{1-u^2}
\end{align*}
So
\[
\int \frac{1}{1-u^2}\,\mathrm{d}u=\operatorname{artanh}u+c
\]

Now let
\[
u=\frac{3x}{2}
\]
Then
\[
x=\frac{2u}{3}, \qquad \mathrm{d}x=\frac{2}{3}\,\mathrm{d}u
\]
Also, when $x=\frac15$, then $u=\frac{3}{10}$, and when $x=\frac25$, then $u=\frac35$.

Therefore
\begin{align*}
\int_{1/5}^{2/5}\frac{1}{4-9x^2}\,\mathrm{d}x
&=\int_{3/10}^{3/5}\frac{1}{4-9\left(\frac{2u}{3}\right)^2}\cdot \frac{2}{3}\,\mathrm{d}u \\
&=\int_{3/10}^{3/5}\frac{1}{4-4u^2}\cdot \frac{2}{3}\,\mathrm{d}u \\
&=\frac{1}{6}\int_{3/10}^{3/5}\frac{1}{1-u^2}\,\mathrm{d}u \\
&=\frac{1}{6}\left[\operatorname{artanh}u\right]_{3/10}^{3/5}
\end{align*}
Using the result from part (i),
\begin{align*}
\int_{1/5}^{2/5}\frac{1}{4-9x^2}\,\mathrm{d}x
&=\frac{1}{6}\left[\frac12\ln\left(\frac{1+u}{1-u}\right)\right]_{3/10}^{3/5} \\
&=\frac{1}{12}\left[\ln\left(\frac{1+u}{1-u}\right)\right]_{3/10}^{3/5} \\
&=\frac{1}{12}\left(\ln\left(\frac{1+\frac35}{1-\frac35}\right)-\ln\left(\frac{1+\frac{3}{10}}{1-\frac{3}{10}}\right)\right) \\
&=\frac{1}{12}\left(\ln 4-\ln\left(\frac{13}{7}\right)\right) \\
&=\frac{1}{12}\ln\left(\frac{28}{13}\right)
\end{align*}
Hence the exact value of the integral is
\[
\frac{1}{12}\ln\left(\frac{28}{13}\right)
\]
\end{enumerate}
\end{tcolorbox}


\newpage

% ---- Question 15 ----
\questionitem
% [Source: _QBT___Solns__Integration_with_Inverse_Trig, Q15]

Show that
\[
\int \frac{x^2+14x+21}{(x-1)(x^2+4x+13)}\,\mathrm{d}x=A\ln|x-1|+B\ln(x^2+4x+13)+D\arctan\left(\frac{x+2}{3}\right)+c
\]
where $A$, $B$ and $D$ are constants to be found. \marks{7} \\

\vspace*{10pt}

\begin{tcolorbox}[colback=white, colframe=scarlet, title={\textbf{Solution}}, breakable, grow to left by=6mm, grow to right by=10mm]
Write
\[
\frac{x^2+14x+21}{(x-1)(x^2+4x+13)}=\frac{P}{x-1}+\frac{Qx+R}{x^2+4x+13}
\]

Then
\[
x^2+14x+21=P(x^2+4x+13)+(Qx+R)(x-1)
\]

Expanding the right-hand side,
\begin{align*}
P(x^2+4x+13)+(Qx+R)(x-1)
&=Px^2+4Px+13P+Qx^2-Qx+Rx-R \\
&=(P+Q)x^2+(4P-Q+R)x+(13P-R)
\end{align*}

Comparing coefficients with $x^2+14x+21$ gives
\begin{align*}
P+Q&=1 \\
4P-Q+R&=14 \\
13P-R&=21
\end{align*}

From the first equation,
\[
Q=1-P
\]

From the third equation,
\[
R=13P-21
\]

Substitute into the second equation:
\begin{align*}
4P-(1-P)+(13P-21)&=14 \\
4P-1+P+13P-21&=14 \\
18P-22&=14 \\
18P&=36 \\
P&=2
\end{align*}

So
\[
Q=1-2=-1,\qquad R=13(2)-21=5
\]

Hence
\[
\frac{x^2+14x+21}{(x-1)(x^2+4x+13)}
=\frac{2}{x-1}+\frac{-x+5}{x^2+4x+13}
\]

Therefore
\[
\int \frac{x^2+14x+21}{(x-1)(x^2+4x+13)}\,\mathrm{d}x
=\int \frac{2}{x-1}\,\mathrm{d}x+\int \frac{-x+5}{x^2+4x+13}\,\mathrm{d}x
\]

Now rewrite the second numerator in terms of the derivative of the quadratic:
\[
-x+5=-\frac12(2x+4)+7
\]

Also,
\[
x^2+4x+13=(x+2)^2+9
\]

So
\begin{align*}
\int \frac{x^2+14x+21}{(x-1)(x^2+4x+13)}\,\mathrm{d}x
&=\int \frac{2}{x-1}\,\mathrm{d}x
-\frac12\int \frac{2x+4}{x^2+4x+13}\,\mathrm{d}x
+7\int \frac{1}{(x+2)^2+9}\,\mathrm{d}x \\
&=2\ln|x-1|-\frac12\ln(x^2+4x+13)+7\int \frac{1}{(x+2)^2+3^2}\,\mathrm{d}x
\end{align*}

Using the standard result
\[
\int \frac{1}{u^2+a^2}\,\mathrm{d}u=\frac{1}{a}\arctan\!\left(\frac{u}{a}\right)+c
\]
with $u=x+2$ and $a=3$,
\[
\int \frac{1}{(x+2)^2+9}\,\mathrm{d}x
=\frac13\arctan\!\left(\frac{x+2}{3}\right)
\]

Hence
\begin{align*}
\int \frac{x^2+14x+21}{(x-1)(x^2+4x+13)}\,\mathrm{d}x
&=2\ln|x-1|-\frac12\ln(x^2+4x+13)+\frac73\arctan\!\left(\frac{x+2}{3}\right)+c
\end{align*}

So the required constants are
\[
A=2,\qquad B=-\frac12,\qquad D=\frac73
\]
\end{tcolorbox}


\newpage

% ---- Question 16 ----
\questionitem
% [Source: _QBT___Solns__Integration_with_Inverse_Trig, Q16]

\begin{figure}[H]
    \centering
\begin{tikzpicture}[scale=1.3]
\def\xmin{-2.2}
\def\xmax{5.2}
\def\ymin{-0.3}
\def\ymax{2.8}
\def\num{9}
\def\peak{1}
\def\den{4}
\def\xL{-1}
\def\xR{4}
\pgfmathsetmacro{\yL}{\num/((\xL-\peak)*(\xL-\peak)+\den)}
\pgfmathsetmacro{\yR}{\num/((\xR-\peak)*(\xR-\peak)+\den)}
\coordinate (O) at (0,0);
\coordinate (L) at (\xL,0);
\coordinate (R) at (\xR,0);
\coordinate (LT) at (\xL,\yL);
\coordinate (RT) at (\xR,\yR);
\draw[thick,-{Stealth[length=3mm]}] (\xmin,0) -- (\xmax,0) node[right] {$x$};
\draw[thick,-{Stealth[length=3mm]}] (0,\ymin) -- (0,\ymax) node[above] {$y$};
\draw[thick] plot[domain=\xL:\xR,samples=200] (\x,{\num/((\x-\peak)*(\x-\peak)+\den)});
\draw[thick] (L) -- (LT);
\draw[thick] (R) -- (RT);
\node[below] at (L) {$-1$};
\node[below] at (R) {$4$};
\node[below left] at (O) {$O$};
\end{tikzpicture}
\end{figure}


The diagram shows the cross-section of a solid wooden beam. \\
Each unit on the axes represents $6$ cm. The curved upper boundary of the beam is modelled by the equation
\[
y=\frac{9}{(x-1)^2+4}
\] \\
The cross-section is bounded by the curve, the $x$-axis and the lines $x=-1$ and $x=4$. \\
Assuming that the beam is a prism of length $240$ cm, find, correct to $3$ significant figures, the volume of the beam. \marks{6} \\

\vspace*{10pt}

\begin{tcolorbox}[colback=white, colframe=scarlet, title={\textbf{Solution}}, breakable, grow to left by=6mm, grow to right by=10mm]
The cross-sectional area is the area under the curve
\[
y=\frac{9}{(x-1)^2+4}
\]
from $x=-1$ to $x=4$.

First find this area in graph units:
\[
A_{\text{units}}=\int_{-1}^{4}\frac{9}{(x-1)^2+4}\,\,\mathrm{d}x
\]

To reduce this to a standard form, let
\[
u=x-1 \qquad \Rightarrow \qquad \mathrm{d}u=\mathrm{d}x
\]
When $x=-1$, $u=-2$, and when $x=4$, $u=3$.

So
\[
A_{\text{units}}=\int_{-2}^{3}\frac{9}{u^2+2^2}\,\,\mathrm{d}u
\]

Using the standard result
\[
\int \frac{1}{a^2+x^2}\,\,\mathrm{d}x=\frac{1}{a}\arctan\!\left(\frac{x}{a}\right)+c
\]
with $a=2$,
\begin{align*}
A_{\text{units}}
&=9\int_{-2}^{3}\frac{1}{u^2+2^2}\,\,\mathrm{d}u \\
&=9\left[\frac{1}{2}\arctan\!\left(\frac{u}{2}\right)\right]_{-2}^{3} \\
&=\frac{9}{2}\left[\arctan\!\left(\frac{u}{2}\right)\right]_{-2}^{3} \\
&=\frac{9}{2}\left(\arctan\!\left(\frac{3}{2}\right)-\arctan(-1)\right) \\
&=\frac{9}{2}\left(\arctan\!\left(\frac{3}{2}\right)+\frac{\pi}{4}\right)
\end{align*}

Now convert to $\text{cm}^2$.

Each unit on each axis represents $6$ cm, so
\[
1\text{ square unit}=6\times 6=36\text{ cm}^2
\]
Hence the cross-sectional area is
\begin{align*}
A
&=36A_{\text{units}} \\
&=36\cdot \frac{9}{2}\left(\arctan\!\left(\frac{3}{2}\right)+\frac{\pi}{4}\right) \\
&=162\left(\arctan\!\left(\frac{3}{2}\right)+\frac{\pi}{4}\right) \\
&\approx 286.447\text{ cm}^2
\end{align*}

The beam is a prism of length $240$ cm, so
\begin{align*}
V&=240A \\
&=240\times 286.447 \\
&\approx 68747.3\text{ cm}^3
\end{align*}

Therefore, correct to $3$ significant figures, the volume of the beam is
\[
6.87\times 10^4\text{ cm}^3
\]
\end{tcolorbox}


\newpage

% ---- Question 17 ----
\questionitem
% [Source: _QBT___Solns__Integration_with_Inverse_Trig, Q17]

Find
\[
\int_1^{e} \frac{1}{x\sqrt{4-(\ln x)^2}} \, \mathrm{d}x
\]
expressing your answer in terms of $\pi$. \marks{8}

\vspace*{10pt}

\begin{tcolorbox}[colback=white, colframe=scarlet, title={\textbf{Solution}}, breakable, grow to left by=6mm, grow to right by=10mm]
Let
\[
u=\ln x
\]
Then
\[
\frac{\mathrm{d}u}{\mathrm{d}x}=\frac{1}{x}
\qquad\Rightarrow\qquad
\mathrm{d}u=\frac{1}{x}\,\mathrm{d}x
\]

We also change the limits:

when \(x=1\), \(u=\ln 1=0\)

when \(x=e\), \(u=\ln e=1\)

So the integral becomes
\[
\int_1^{e} \frac{1}{x\sqrt{4-(\ln x)^2}} \,\mathrm{d}x
=
\int_0^{1} \frac{1}{\sqrt{4-u^2}} \,\mathrm{d}u
\]

Now use the standard result
\[
\int \frac{1}{\sqrt{a^2-u^2}} \,\mathrm{d}u
=
\arcsin\!\left(\frac{u}{a}\right)+c
\]

Here \(a=2\), so
\[
\int_0^{1} \frac{1}{\sqrt{4-u^2}} \,\mathrm{d}u
=
\left[\arcsin\!\left(\frac{u}{2}\right)\right]_0^1
\]

Therefore
\[
\left[\arcsin\!\left(\frac{u}{2}\right)\right]_0^1
=
\arcsin\!\left(\frac12\right)-\arcsin(0)
=
\frac{\pi}{6}-0
=
\frac{\pi}{6}
\]

Hence the value of the integral is
\[
\frac{\pi}{6}
\]
\end{tcolorbox}


\end{enumerate}
\end{document}
